%-*-latex-*-
%-*-latex-*-
\newcommand\COURSE{ciss350}
\newcommand\ASSESSMENT{q10901}
\newcommand\ASSESSMENTTYPE{Quiz}
\newcommand\POINTS{\textwhite{xxx/xxx}}

\input{myquizpreamble}
\input{yliow}
\input{\COURSE}
%-*-latex-*-

%-*-latex-*-

%-*-latex-*-
\renewcommand\TITLE{\ASSESSMENTTYPE \ \ASSESSMENT}

\renewcommand\EMAIL{}
\renewcommand\AUTHOR{eagraham1@cougars.ccis.edu}

\textwidth=6in
\begin{document}
\topmatterthree

Here's how you describe (in text) for a binary tree with a heap tree structure.
For this tree
\begin{console}[fontsize=\footnotesize]
              (10)
           /        \
        (6)          (7)
       /   \        /   \
    (5)     (2)  (0)     (3)
   /
(1)
\end{console}
you write
\begin{console}[fontsize=\footnotesize]
{10, 6, 7, 5, 2, 0, 3, 1}
\end{console}
i.e., list the values using breadth-first traversal (like what we did in class).

%------------------------------------------------------------------------------
\nextq
Describe (using text) the following maxheap after you have inserted 9 to it.
\begin{console}[fontsize=\footnotesize]
              (10)
           /        \
        (6)          (7)
       /   \        /   \
    (5)     (2)  (0)     (3)
   /
(1)
\end{console}

\ANSWER
\begin{answercode}
{10, 6, 7, 5, 2, 0, 3, 1, 9}
{10, 6, 7, 9, 2, 0, 3, 1, 5}
{10, 9, 7, 6, 2, 0, 3, 1, 5}
\end{answercode}

%------------------------------------------------------------------------------
\nextq
Describe (using text) the following \textit{min}heap after you have inserted 0
to it.
\begin{console}[fontsize=\footnotesize]
              (1)
           /       \
        (3)         (2)
       /   \       /   \
    (6)     (4) (8)     (5)
   /   \
(9)     (7)
\end{console}

\ANSWER
\begin{answercode}
[1, 3, 2, 6, 4, 8, 5, 9, 7, 0]
[1, 3, 2, 6, 0, 8, 5, 9, 7, 4]
[1, 0, 2, 6, 3, 8, 5, 9, 7, 4]
[0, 1, 2, 6, 3, 8, 5, 9, 7, 4]
\end{answercode}

%------------------------------------------------------------------------------
\nextq
Describe (using text) the following maxheap after you have performed
extract-root on it.
\begin{console}[fontsize=\footnotesize]
              (10)
           /        \
        (6)          (7)
       /   \        /   \
    (5)     (2)  (0)     (3)
   /
(1)
\end{console}

\ANSWER
\begin{answercode}
[10, 6, 7, 5, 2, 0, 3, 1]

[1, 6, 7, 5, 2, 0, 3]
[7, 6, 1, 5, 2, 0, 3]
[7, 6, 3, 5, 2, 0, 1]
\end{answercode}

%------------------------------------------------------------------------------
\nextq
Complete the following maxheap insert function where \verb!std::vector< int >!
is used for the array representation of the maxheap.

\ANSWER
\begin{answercode}
#include <iostream>
#include <string>
#include <vector>

std::ostream & operator<<(std::ostream & cout, const std::vector< int > & v)
{
    cout << '[';
    std::string delim = "";
    for (auto & e: v)
    {
        cout << delim << e; delim = ", ";
    }
    cout << ']';
    return cout;
}

int PARENT(int i)
{
  return (i - 1) / 2;
}

int LEFT(int i)
{
  return 2 * i + 1;
}

int RIGHT(int i)
{
  return 2 * i + 1;
}

void maxheap_insert(std::vector< int > & v, int x)
{
  v.push_back(x);

  int n = v.size() - 1;
  while(n > 0 && v[n] > v[PARENT(n)]
  {
    int temp = v[n];
    v[n] = v[PARENT(n)];
    v[PARENT(n)];
    n = PARENT(i);
  }
}

int main()
{
    std::vector< int > heap;
    std::vector< int > data {5, 3, 1, 2, 4, 6, 0};
    for (auto & e: data)
    {
        maxheap_insert(heap, e);
        std::cout << heap << '\n';
    }
    return 0;
}
\end{answercode}

%------------------------------------------------------------------------------
%-*-latex-*-
\newpage
\input{instructions}
\end{document}

